\documentclass[../../../include/open-logic-section]{subfiles}

\begin{document}

\olfileid[tr]{sth}{choice}{banach}
\olsection{Banach--Tarski Paradoksu}

Boolos'un Yerine Koyma'yı gerekçelendirmeye çalıştığı gibi biz de
\emph{dışsal} değerlendirmelere başvurarak Seçim'i gerekçelendirmeyi
deneyebiliriz (bkz. \olref[replacement][extrinsic]{sec}). Sonuçta Seçim'i
benimsemenin birçok istenir sonucu vardır: her kardinal sayıyı
karşılaştırabilmek, her kümeyi iyi sıralayabilmek, kardinal sayıları belirli
bir sıra sayısı türü olarak ele alabilmek vb.

Fakat bazen Seçim'in \emph{istenmeyen} sonuçları olduğu öne sürülür.
Bu çoğunlukla \cite{BanachTarski1924} sonucundan kaynaklanır.

\begin{thm}[Banach--Tarski Paradoksu ($\ZFC$ içinde)]
Her top sonlu sayıda parçaya ayrılabilir ve bu parçalar (döndürme ve
öteleme yoluyla) yeniden bir araya getirilerek o topun iki kopyası
oluşturulabilir.
\end{thm}
\noindent
İlk bakışta bu oldukça şaşırtıcıdır. İki topun hacmi açıkça özgün topun
hacminin \emph{iki katıdır}. Fakat rijit hareketler---döndürme ve
öteleme---hacmi değiştirmez. Dolayısıyla Banach--Tarski yoktan yeni madde
var edebilmemize izin veriyor görünür.

Durum daha da kötüleşir.\footnote{Bkz.
\citet[Theorem 3.12]{Wagon2016}.} Benzer bir uslamlama, bir bezelyenin sonlu
sayıda parçaya ayrılabileceğini ve bu parçaların (döndürme ve öteleme
yoluyla) yeniden bir araya getirilerek Big Ben biçiminde ve büyüklüğünde
bir varlık oluşturulabileceğini gösterir.

Ne var ki bunların hiçbiri tek başına $\ZF$ içinde geçerli değildir.
\footnote{Banach--Tarski, Seçim'den kesinlikle daha zayıf ilkelerle de
ispatlanabilir; bkz. \citet[303]{Wagon2016}.} Öyleyse bir kararla karşı
karşıyayız: Seçim'i reddetmek ya da ``paradoks''la yaşamayı öğrenmek.

``Paradoks''la yaşamayı öğrenmemiz gerektiğini öne süreceğiz. Gerçekten
bunun büyük bir paradoks olduğunu düşünmüyoruz. Özellikle aşağıdaki
sonuçların herhangi birinden niçin daha fazla ya da daha az paradoksal
olduğunu göremiyoruz:\footnote{\citet[276--7]{Potter2004},
\citet[16]{Weston2003} ve \citet[31, 308--9]{Wagon2016} başka örnekler
kullanarak benzer noktalar ileri sürer.}
\begin{enumerate}
	\item $(0,1)$ aralığında $\Real$'deki kadar nokta vardır.
	\\\emph{İspat}: $\tan(\pi(r-\nicefrac{1}{2}))$ fonksiyonunu ele alın.
	\item Bir doğru üzerinde bir karedeki kadar nokta vardır.
	\\Bkz. \olref[his][set][pathology]{sec} ve
	\olref[his][set][cantorplane]{sec}.
	\item Düzlem dolduran eğriler vardır.
	\\Bkz. \olref[his][set][pathology]{sec} ve
	\olref[his][set][hilbertcurve]{sec}.
\end{enumerate}
Bu üç sonucun hiçbiri Seçim'i gerektirmez. Gerçekten artık bunları yalnızca
şaşırtıcı ve güzel matematik parçaları olarak görüyoruz. Belki
Banach--Tarski Paradoksu'na karşı da aynı tutumu benimsemeliyiz.

Kuşkusuz burada teknik bir gözlem gerekir; fakat bu yalnızca soğukkanlı
olmayı gerektirir. Rijit hareketler hacmi korur. Dolayısıyla topun ayrıldığı
beş\footnote{Paradoksu ``sonlu sayıda parça'' cinsinden ifade ettik.
Gerçekte \citet{Robinson1947}, ayrıştırmanın \emph{beş} parçayla (fakat
daha azıyla değil) yapılabileceğini ispatladı. Bir ispat için bkz.
\citet[pp.~66--7]{Wagon2016}.} parçanın hepsi \emph{ölçülebilir} olamaz.
Kabaca söylersek bu ayrı ayrı parçalara bir hacim atamak anlamsızdır.
Bunları gözde canlandırılamayan, noktaların ``sonsuz saçılımları'' olarak
düşünmelisiniz. Böyle ``sonsuz saçılmış'' kümeleri tasarlamak belki
``tuhaf''tır. Fakat varlıkları, önceden kullanılabilir kümelerin
\emph{bütün olanaklı} koleksiyonlarını oluşturmamızı söyleyen ve
\stagesacc{} içinde somutlaşan buyruktan çıkıyor görünür.

Bunların hiçbiri ikna edici gelmediyse Banach--Tarski Paradoksu'nu
benimsemek lehine son bir (dışsal) sav şudur: bütün zamanların en iyi
matematik şakasını doğrudan gerektirir:
\begin{enumerate}
	\item[] \emph{Soru}. ``Banach--Tarski''nin bir anagramı nedir?
	\item[] \emph{Yanıt}. ``Banach--Tarski Banach--Tarski''.
\end{enumerate}

\end{document}
